% from ega4-14-fr.tex
\oldpage[IV-3]{222}

ouvert dans $Y$ et contient $y$, donc est égal à $Y$. On peut ainsi remplacer $X$
par $U$ pour prouver la proposition.

\emph{V) Fin de la démonstration.} -- Nous supposons donc $Y$ local, intègre et
géométriquement unibranche, $X$ intègre et affine, $f$ quasi-fini, surjectif,
birationnel et universellement ouvert~; il en résulte que $f$ est automatiquement
\emph{séparé}. En vertu du Main theorem (8.12.6), il existe une factorisation
$X\xrightarrow{j}Z\xrightarrow{u}Y$, où $j$ est une immersion ouverte et $u$ un
morphisme \emph{fini}. Remplaçant en outre $Z$ par l'image fermée de $j$
(\textbf{I}, 9.5.10), on peut supposer que $Z$ est \emph{intègre}~; comme $u$ est
fini et birationnel et $Y$ géométriquement unibranche, il résulte de
(\textbf{III}, 4.3.5 et 4.3.4) que $u$ (et par suite $f$) est un morphisme
\emph{radiciel}~; comme $f$ est universellement ouvert et surjectif, c'est donc un
\emph{homéomorphisme universel}. Par suite (2.4.5, (ii)) $f$ est un morphisme
\emph{fini}. Mais alors on répond aux conditions de l'énoncé avec $Y'$ \emph{intègre}
en prenant simplement $Y'=X$ et $g=f$, la $Y'$-section étant le morphisme diagonal
$\Delta_f$. C.Q.F.D.

Ce résultat peut être utilisé pour développer des critères de «~descente~» de
diverses propriétés par les morphismes universellement ouverts et surjectifs.
Signalons en particulier le critère suivant, dû à D.~Mumford~:

\begin{corollary}[14.5.11]
\EGAReaderTarget{ega4.c04.cor.14-5-11}{IV.14.5.11-fr}
Soient $Y$ un préschéma localement noethérien, $f:X\to Y$ un morphisme localement
de type fini, $g:Y'\to Y$ un morphisme localement de type fini, universellement
ouvert et surjectif~; posons $X'=X\times_Y Y'$, $f'=f_{(Y')}:X'\to Y'$. Alors,
pour que $f$ soit affine, il faut et il suffit que $f'$ le soit.
\end{corollary}

On doit seulement prouver que la condition est suffisante. La question étant locale
sur $Y$, on peut supposer $Y$ affine, donc noethérien. En vertu de (14.5.10), il
existe un morphisme fini surjectif $h:Y_1\to Y$ tel que, si l'on pose
$Y'_1=Y'\times_Y Y_1$, et $g'=g_{(Y_1)}:Y'_1\to Y_1$, tout point de $Y_1$ admette
un voisinage ouvert $U_1$ tel qu'il existe une $U_1$-section de $g'^{-1}(U_1)$. Si
l'on pose $X_1=X\times_Y Y_1$, la projection canonique $p:X_1\to X$ est un
morphisme fini surjectif, donc, si l'on prouve que $X_1$ est un schéma affine, il en
résultera d'abord que $X$ est quasi-compact, donc noethérien, puis que $X$ est affine
en vertu du théorème de Chevalley (\textbf{II}, 6.7.1). Il suffit par suite de
prouver que le morphisme $f_1=f_{(Y_1)}:X_1\to Y_1$ est affine. Or, si l'on pose
\[
  X'_1=X_1\times_{Y_1}Y'_1=X'\times_{Y'}Y'_1
  \qquad\text{et}\qquad
  f'_1=(f_1)_{(Y'_1)}=f'_{(Y'_1)}:X'_1\to Y'_1,
\]
$f'_1$ est affine en vertu de l'hypothèse. On est donc ramené à prouver le corollaire
en y remplaçant $Y$, $X$, $f$ et $Y'$ par $Y_1$, $X_1$, $f_1$ et $Y'_1$,
autrement dit, il suffit de prouver que $f$ est affine lorsqu'on fait de plus, dans
l'énoncé de (14.5.11), l'hypothèse que tout point de $Y$ admet un voisinage ouvert
$U$ tel qu'il existe une $U$-section de $g^{-1}(U)$. La question étant locale sur
$Y$, on peut même supposer qu'il existe une $Y$-section $s$ de $Y'$. Or, on a le
lemme élémentaire suivant (valable dans toute catégorie admettant des produits
fibrés)~:

\begin{lemma}[14.5.11.1]
\label{IV.14.5.11.1-fr}
Soient $f:X\to S$, $g:Y\to S$ deux morphismes, $p_1:X\times_S Y\to X$,
$p_2:X\times_S Y\to Y$ les projections canoniques. Si $s:S\to Y$ est une section
de $g$, alors $s'=(1,s\circ f)_S$ est une section de $p_1$ et $X$, muni des
morphismes $f:X\to S$ et $s':X\to X\times_S Y$, s'identifie au produit des
$Y$-préschémas $S$ et $X\times_S Y$ pour les morphismes $s:S\to Y$ et
$p_2:X\times_S Y\to Y$.
\end{lemma}


% (continued in ega4-15-fr.tex, before its first page marker)
\section{Étude des fibres d'un morphisme universellement ouvert}
\label{section:IV.15-fr}

\subsection{Multiplicités des fibres d'un morphisme universellement ouvert}
\label{subsection:IV.15.1-fr}

\begin{proposition}[15.1.1]
\label{IV.15.1.1-fr}
